\section{VRAM Layouts: Decoupling the Disk Rate from the Memory Rate}
\label{sec:layouts}

The rate deciding which models fit on a GPU is that of the format the kernel
reads, not the artifact's on-disk rate. We map it: four byte formats of ours read by
five kernels (one gets two decoders), two deployed competitors at 4 and 2
bits, FP16 and cuBLAS controls, a floor kernel reading no weights, one
harness, one byte-accounting convention.

\textbf{Methodology.} All arms run in one process on an NVIDIA L40S,
interleaved each round in a fixed dispatch order; 7 rounds, the first 2
discarded. Each speedup is the \emph{median of per-round ratios} with its
range, never a quotient of two minima --- rounds that never coexisted. All kernels use zero local memory; every row of every arm is
verified against f64 before timing (\S\ref{sec:decoder}).

Byte accounting is the kernel's: what each arm streams per matrix (payload,
exception side-channels, full-precision tail columns, per-row scales). Three
asymmetries follow, uncorrected. Per-class tables are excluded as resident
constants, kilobytes against gigabytes of stream, though \textsc{Golay70}
loads two. The FP16 control bills its tail at f16 where ours bill
theirs at f32, and the 4-bit competitor bills none: both understate our arms.
Bandwidth follows the bench's convention, each arm's fastest kept round; on
medians every percentage is unchanged.

Adding an arm perturbs the others, so the run has two phases: nine arms, then
those nine plus the tenth. Across phases the nine common ratios move by at
most 0.36\%, this table's resolution; every difference argued from it
stands at several times that grain. \emph{Every row comes from that one
run}: appending the 2-bit competitor to the earlier seven-arm table would set
a row from one process beside rows from another, the comparison this
methodology forbids, so we re-measured with every arm present. The
incumbents reproduce within dispersion ten days later (\textsc{Planes14}
$2.15\times$ both times, hoisted \textsc{Golay70} $1.77\times$ then
$1.78\times$). Data: \texttt{docs/data/echelle-formats.csv}; raw log
\texttt{docs/mesures/f2-p3-qtip-banc-2026-08-21.txt}.
$^{\dagger}$The bench prints ratios against FP16, not another arm, so the one
cross-arm ratio we state (\S\ref{sec:qtip}) uses the two medians with its
range bounded outward by both arms' min--max --- the only ratio here not
formed round by round, widened conservatively.

\begin{table}[t]
\centering
\small
\setlength{\tabcolsep}{3pt}
\begin{tabular}{lrrrrrrl}
\toprule
Kernel & b/weight & GB & med.\ ms & GB/s & of bound & vs FP16 & status \\
\midrule
Floor (no weights read) & 0.159 & 0.07 & 2.306 & 31 & 5\% & 4.77$\times$ [4.76--4.77] & launch geometry \\
FP16 (control) & 16.000 & 7.27 & 10.994 & 661 & --- & 1.00$\times$ & \\
FP16 via cuBLAS & 16.000 & 7.27 & 10.830 & 672 & 102\% & 1.02$\times$ [1.02--1.02] & checks the control \\
\midrule
\textsc{Slot32} & 5.510 & 2.50 & 5.820 & 431 & 65\% & 1.89$\times$ [1.89--1.89] & first working \\
\textsc{Planes14} & 4.804 & 2.18 & 5.103 & 428 & 65\% & \textbf{2.15$\times$} [2.15--2.16] & \textbf{in production} \\
\textsc{Planes12x} & 4.342 & 1.97 & 5.492 & 359 & 54\% & 2.00$\times$ [2.00--2.00] & best exact rate \\
\textsc{Golay70} & 3.589 & 1.63 & 8.187 & 199 & 30\% & 1.34$\times$ [1.34--1.34] & discarded \\
\textsc{Golay70}, hoisted & 3.589 & 1.63 & 6.182 & 264 & 40\% & 1.78$\times$ [1.77--1.78] & still discarded \\
\midrule
AWQ w4g128 & 4.179 & 1.90 & 3.252 & 584 & \textbf{88\%} & 3.38$\times$ [3.37--3.38] & 4-bit competitor \\
QTIP 2-bit & 2.000 & 0.91 & \textbf{2.246} & 405 & 61\% & 4.89$\times$ [4.89--4.90] & \textbf{2-bit competitor} \\
\bottomrule
\end{tabular}
\caption{The layout scale with both competitors: one token of the 252 model
projections, single L40S, ten arms in one process, median of
per-round ratios with ranges. \emph{Of bound} is achieved bandwidth as a
fraction of the 661~GB/s the FP16 control reaches on these shapes --- the
comparable quantity, a ratio against FP16 being mechanically higher for the
arm reading least. \emph{The control is our own FP16 matvec, checked}:
cuBLAS on the same weights and rounds runs 1.5--2.4\% faster, leaving every
ratio standing; ours stays the anchor. The floor kernel runs the same 252
launches over the same shapes and reads no weight bytes: its $4.77\times$ is
our launch geometry's cost before any format is chosen,
and the 2-bit competitor, in its own grid, beats it. Kernel accounting;
whole-model bits per parameter (\S\ref{sec:evaluation}) are a different
denominator.}
\label{tab:layouts}
\end{table}

\begin{figure}[t]
\centering
\includegraphics[width=0.82\linewidth]{fig1_layout_scale.pdf}
\caption{Speedup over FP16 against in-VRAM rate; error bars min--max over
kept rounds. The dashed hyperbola is the ceiling: what an arm at that rate
would reach at the FP16 control's bandwidth on the same shapes. An arm's
fraction of its own byte bound is its height as a \emph{proportion} of that
curve, hence the log axis: \textsc{Slot32} and \textsc{Planes14}, both at
65\%, sit equally far below the ceiling at very different rates. The arrow
is the second attack on \textsc{Golay70}: the per-slot coset decode hoisted
to a per-block prologue at zero change to the stored bytes, 199 to 264~GB/s,
still short. The horizontal rule is the floor: our launch geometry
over the same 252 launches with no weights to read. Every arm of ours lies
below it, as it must; the 2-bit competitor lies \emph{above} it
(\S\ref{sec:qtip}) --- kernel structure, not card.}
\label{fig:layouts}
\end{figure}

\textbf{\textsc{Slot32}} stores decoded fields at fixed offsets in a
bit-packed, byte-aligned record of $9 + g + 24L$ bits: 9-bit class
identifier, $g=1$ gain bit, 24-bit sign mask, one 24-bit one-hot mask per
non-zero magnitude level --- 130 bits at the widest class ($L=5$), 106 at
$L=4$. The kernel addresses a record by byte offset and reads a fixed
five-word (20~B) window, an alignment the format guarantees rather than
checks. Simplest divergence-free format, our first to beat FP16
($1.89\times$), but at 5.510~b/weight \emph{more} memory than a 4-bit
format: fast, pointless at scale.

\textbf{\textsc{Planes14}} removes the waste in the one-hot level masks
$m_1..m_4$: blocks carry three to five magnitude levels, so three
binary bit planes name each coordinate's level where one-hot spent up to four
24-bit masks. With a uniform 14-byte stride, no level cap and no per-group
base pointers, the block shrinks to 4.804~b/weight and, the kernel being
memory-bound, runs \emph{faster}: $1.14\times$ over \textsc{Slot32} at
bit-identical decoded content ($1.14\times$ in five independent jobs, union
of ranges 1.13--1.15) and $2.15\times$ over FP16. Achieved bandwidth is
unchanged (431 to 428~GB/s): time falls as bytes do, and two access
geometries at one bandwidth is evidence that record geometry is not what
binds here (\S\ref{sec:attribution}).

\textbf{\textsc{Planes12x}} adds a sparse overlay: blocks within a stricter
level budget take a 12-byte fast path, rare violators an exact side-channel
correction. Reconstruction is bit-exact, and the alternative is not free:
\emph{hard-capping} the level count at four (a plain two-plane 12-byte format
at about 4.1~b/weight) was measured to cost $+4.75\%$ perplexity, enough to
fall behind QTIP. The overlay gets 4.342~b/weight at $2.00\times$ with zero
quality change.

\textbf{\textsc{Golay70}: a negative result, attacked twice.} The sign mask
is redundant --- signs on the codeword support are determined up to the Golay
coset --- so 12 code bits in place of 24 sign bits yield 3.589~b/weight,
reconstruction-exact on all 150{,}681{,}600 blocks of the artifact. But
recovering them resolves the coset \emph{per slot} and the kernel stops being
memory-bound: 199~GB/s, 30\% of its byte bound, $1.34\times$, below the
$1.6\times$ criterion fixed before measuring in the commit message landing
the layout\ifanon{ (\texttt{caef2ac})}{} --- antecedence resting on a commit
date, where the replacement criterion below carries a trusted timestamp.

We attacked the identified cause. The two cosets differ by a function of the
block's three 24-bit words, so it hoists into a per-block prologue of about a
dozen integer operations, leaving three mask tests and a negation per slot; \emph{not one stored byte changes}, the identity proved
slot by slot and block by block. The repair works and is not enough: 199 to
264~GB/s, 30\% to 40\%, $1.34\times$ to $1.78\times$. Halving the
per-slot decode returned 2.01~ms of the 5.72~ms by which the layout missed
its DRAM floor: right diagnosis, bar still missed. The arms stand above their
own byte floors by 1.81, 2.51 and 3.72~ms (\textsc{Planes14},
\textsc{Planes12x}, hoisted \textsc{Golay70}), the excess over
\textsc{Planes14}'s gap growing from 0.71~ms at 3.4\% exceptions to 1.91 at
7.4\%: faster than in proportion to a correction pass that lies outside the
path this repair touched. The design space is not empty; the layout missed a
pre-registered bar twice.

The bar had moved between attempts: with a deployed 4-bit kernel faster and
smaller than our production layout (below), a speed threshold could no longer
settle a memory argument. The replacement
($\geq 2.0\times$, the speed of \textsc{Planes12x}, plus a $\geq 20\%$
whole-model memory margin) was timestamped before measuring. The memory
condition held by arithmetic; the speed condition missed, $1.77\times$
against $2.0\times$ in the deciding campaign (re-measured at $1.78\times$ ten
days later: dispersion, not a second verdict). Our $1.9$--$2.4\times$ estimate, from an instruction count, was
optimistic for the reason flagged then as its weakest: integer and float
issue share ports.

The scale is nonlinear: \emph{bits saved are speed gained on a memory-bound
kernel, until the decode saving them stops being shifts and masks} --- for
lattice codes here, between 4.8 and 4.3~b/weight.

\textbf{Where a deployed 4-bit kernel sits.} The first of
Table~\ref{tab:layouts}'s two outside arms is the AWQ w4g128 GEMV
kernel~\citep{awq2024}, ported into our harness unchanged (own grid, own
binary16 output), timed in the same rounds and process as ours. It
reads 4.179~b/weight, less than three of our four byte formats, so a ratio
against FP16 flatters it by construction and we do not lean on the
$3.38\times$; against the two \textsc{Golay70} rows, which read less, the
same effect runs the other way and they lose regardless. The comparable quantity
is the fraction of the byte bound, this section's result: \textbf{88\% for
the 4-bit kernel against 65\% for our best layout}.

Two asymmetries qualify that number. The AWQ payload is \emph{synthesized}
(our own weights requantized to w4g128): sound for timing --- no
data-dependent branch, traffic depending only on the shapes --- and unsound
for quality, which this arm claims none of. w4g128 also quantizes every
column, carrying no tail where our rates do; billing \textsc{Planes14}
without tail and row scales gives 63\% of bound rather than 65\%, widening
the gap to 25 points. We keep 65\%: what the shipped format streams.

\subsection{And where the 2-bit state of the art sits}
\label{sec:qtip}

Our motivation is an inherited sentence --- the original LLVQ kernel ``is
reported slower than QTIP's'' --- a citation, not a measurement, which
Table~\ref{tab:layouts}'s last row closes. We fetch the published QTIP
inference kernel at a pinned commit (GPL~v3, hash-verified, measured but not
redistributed), give it its own launch geometry as shipped, tune nothing either way, and time it in the same rounds and process as the rest. Its payload is pseudo-random: again sound for timing (no
data-dependent branch) and unsound for quality, on which this arm says
nothing. All $1\,105\,920$ output rows are checked
against the f64 reference at \emph{our} tolerance, not the looser one we
concede to arms emitting binary16: worst error $5.4\times10^{-8}$ against
$10^{-5}$.

\textbf{The inherited claim holds, against us.} Formed as every other
comparison here --- same phase, process and shapes --- the QTIP kernel runs
$\mathbf{2.27\times}$ [2.27--2.28] faster than our production
layout.$^{\dagger}$ The three possible verdicts were written and timestamped
before measuring; this is the one saying our motivation was right about the
world, wrong about us.

\textbf{The mechanism is neither quality nor the format's rate.} Our artifact
stores 2.07 bits per weight on disk (\S\ref{sec:evaluation}) and QTIP's
2.000 --- the same operating point as \emph{codes}, within noise on
perplexity at 4B (Appendix~\ref{app:scale}). A $\Lambda_{24}$ index is a combinatorial address --- class, Golay
codeword, mixed-radix rank, sign rules --- which this section twice showed does not
decode at memory speed (\textsc{Golay70}, $1.34\times$ then
$1.78\times$ against $2.0\times$). The gap opens at load time: the served
path \emph{unfolds} it into bit planes read with shifts and masks, paying 4.804
b/weight in VRAM. QTIP never unfolds: its state is a sliding 16-bit window
over the bitstream, consecutive states overlapping by 12 bits, only $kV = 4$
new bits stored, a weight decoding in a few instructions and one lookup in a
2~KiB table. It reads 0.91~GB where we read 2.18 on the same matrices.
\emph{The Leech lattice packs as well as the trellis and unfolds
worse} --- not a quality deficit, not an implementation detail.

\textbf{It also breaks our own floor, the more useful result.}
Table~\ref{tab:layouts}'s first row is a kernel of ours running the same 252
launches over the same shapes reading \emph{no weight bytes at all}:
2.306~ms, $4.77\times$ FP16, our assumed ceiling on what any format could
buy. The 2-bit competitor does the same projections, reading
0.91~GB, in 2.246~ms --- faster than our kernel that reads nothing, by 2.7\%
against a 0.36\% resolution. The floor was never a property of the machine
but of our launch geometry --- one warp per output row across 252 launches
--- and a differently shaped kernel walks through it. That reframes
\S\ref{sec:attribution}: its largest term, unmasked latency and occupancy at
39\% of the gap, is not a cost of lattice decoding, and the fact that it has
never been attacked becomes a measured opportunity rather than an estimated
one.

\textbf{And it marks the edge of our metric.} The fraction of
the byte bound compares an arm to $16/b$ times the control, interpretable
only while $16/b$ stays under the floor, i.e.\ while $b > 16/4.77 =
3.35$~b/weight. Every arm here had been above it; at 2.000, QTIP is the first
below, its byte bound $8.00\times$ against a floor of $4.77\times$. We
pre-registered the prediction that its fraction could not exceed 59.6\%,
structurally, whatever the kernel; it measures 61\%. Both halves matter: the metric does
saturate as predicted, and the prediction was still wrong, resting on a floor
we did not yet know was ours, not the card's. Ranking the last two rows down the \emph{of bound} column would be
true and useless: where byte bounds are comparable it is the right
instrument, but at 2 bits the honest quantity is the time.

\subsection{How close to the floor? A measured attribution}
\label{sec:attribution}

The instrumented job carries its own FP16 control (7.27~GB at 662~GB/s, a
rounding step from Table~\ref{tab:layouts}'s 661: a floor from one process
and a time from another is the subtraction this methodology forbids). Against
it the 2.50~GB \textsc{Slot32} reads put its DRAM floor at 3.78~ms, while the
instrumented arm takes 5.82~ms (5.820 in Table~\ref{tab:layouts}:
instrumentation costs nothing measurable). That \textbf{2.04~ms} gap is the
kernel's alone; three ``floor'' kernels, each adding one stage, attribute
it:

\begin{table}[h]
\centering
\small
\begin{tabular}{lrr}
\toprule
Stage & $\Delta$ ms & share of gap \\
\midrule
unmasked latency / occupancy & 0.803 & 39\% \\
payload streaming (group bases $+$ five-word window) & 0.681 & 33\% \\
residual decode & ${\sim}0.396$ & 19\% \\
24 shared-memory reads of $x$ (bank conflicts) & 0.118 & 6\% \\
class-table gather & 0.041 & 2\% \\
\bottomrule
\end{tabular}
\caption{Attribution of the 2.04~ms between \textsc{Slot32} and its DRAM
floor. \emph{A decomposition by subtraction of floor kernels, not a
hardware profile}; the residual-decode line is by difference. CUDA
events bound the \emph{host} share at 0.1--0.2\% of the wall on every arm:
the latency term is not the host failing to feed the device but gaps between
kernels on the stream. A hardware profile would settle the rest and could
not be obtained: counters inaccessible on our compute platform
(\texttt{ERR\_NVGPUCTRPERM}), Nsight Compute installed and attached. Data: \texttt{docs/mesures/attribution-cuda-2026-08-05.txt},
\texttt{docs/\allowbreak mesures/\allowbreak fusion-qkv-cuda-2026-08-05.txt}
and \texttt{docs/\allowbreak mesures/\allowbreak f3-events-2026-08-19.txt}.}
\label{tab:attribution}
\end{table}

A source-level instruction count on the same kernel disagrees on every line:
bank conflicts at 0.3--2.9~ms (they are 0.118), the table gather
overestimated twelvefold, arithmetic ``out of the picture'' when decode and
the latency it fails to hide are the largest recoverable terms.

One term is already recovered: fusing $q/k/v$ and
$\textit{gate}/\textit{up}$ into row-concatenated launches takes the model
from 252 kernels to 144 and returns \textbf{0.803~ms}, 39\% of the gap,
output bit-identical on $921\,600$ rows, bytes read identical: a purely
geometric gain. We publish no cross-arm ratio from it: the FP16 arm has no
fused counterpart, and the fused kernel is \textsc{Slot32}-only, off the
served path of \S\ref{sec:integration}.

The $2.15\times$ of Table~\ref{tab:layouts} is therefore not a format
ceiling: measured on \textsc{Slot32}, its largest term is launch geometry,
orthogonal to the record format, and roughly a third of the gap has never
been attacked. The 4-bit competitor's 88\% says the same from outside ---
much of the 23 points between us is not lattice decoding. The 2-bit
competitor settles it from further out still: running the same 252
projections in its own grid, it finishes ahead of our kernel that reads no
weights (\S\ref{sec:qtip}), so that 39\% is recoverable by construction and
not by a format.

\subsection{The same layouts on a second memory hierarchy}
\label{sec:a100}

Everything above is one card, and a byte saved becomes time saved only if the
kernel is bandwidth-bound: \S\ref{sec:layouts}'s scale is a claim about a
\emph{pair}, format and memory hierarchy. We re-ran nine of these
arms, from the same binary and artifact, on an NVIDIA A100 (HBM2e,
\texttt{sm\_80}) against the L40S of Table~\ref{tab:layouts} (GDDR6,
\texttt{sm\_89}). Kernels are compiled for the target at startup; nothing in
their text changed, and line-by-line f64 verification returns the
\emph{same} worst errors on both cards, so only the machine differs.

\begin{table}[t]
\centering
\small
\begin{tabular}{lrrrrr}
\toprule
 & \multicolumn{2}{c}{L40S (GDDR6)} & \multicolumn{2}{c}{A100 (HBM2e)} & \\
\cmidrule(lr){2-3}\cmidrule(lr){4-5}
Arm & GB/s & vs.\ FP16 & GB/s & vs.\ FP16 & b/weight \\
\midrule
Floor (no weights read) &  31 & 4.77$\times$ &   18 & 1.68$\times$ & 0.159 \\
FP16 (control)          & 661 & 1.00$\times$ & 1052 & 1.00$\times$ & 16.000 \\
FP16 via cuBLAS         & 672 & 1.02$\times$ & 1204 & 1.14$\times$ & 16.000 \\
AWQ w4g128 (competitor) & 584 & 3.38$\times$ &  501 & 1.82$\times$ & 4.179 \\
\textsc{Slot32}         & 431 & 1.89$\times$ &  266 & 0.73$\times$ & 5.510 \\
\textsc{Planes14}       & 428 & 2.15$\times$ &  250 & 0.79$\times$ & 4.804 \\
\textsc{Planes12x}      & 359 & 2.00$\times$ &  209 & 0.73$\times$ & 4.342 \\
\textsc{Golay70} v2     & 264 & 1.78$\times$ &  147 & 0.62$\times$ & 3.589 \\
\textsc{Golay70} v1     & 199 & 1.34$\times$ &  104 & 0.44$\times$ & 3.589 \\
\bottomrule
\end{tabular}
\caption{Same arms, two memory hierarchies, same binary, artifact and
protocol. \emph{The two \textsc{vs.\ FP16} columns are not
comparable and never divided}: each is formed round by
round against its own card's control. Data:
\texttt{docs/data/echelle-formats.csv} and
\texttt{docs/\allowbreak data/\allowbreak echelle-formats-a100.csv};
journal \texttt{docs/mesures/f4-a100-2026-08-18.txt}. What transfers is the
\emph{ordering}, which is all the text relies on. The 2-bit competitor has no
A100 point and is absent from both columns rather than half-filled.}
\label{tab:a100}
\end{table}

\textbf{The advantage does not transfer, and GB/s shows why.} On
the A100 no decoding arm reaches the FP16 baseline: our production layout
runs at $0.79\times$ where it ran at $2.15\times$ --- not a faster baseline
outrunning them, but decoders slower \emph{in absolute terms} on the card
with more bandwidth: \textsc{Planes14} falls from 427 to 250~GB/s and every
LLVQ arm with it, while FP16 rises from 661 to 1052 and cuBLAS to 1204. A
memory-bound kernel cannot behave that way. These arms are compute-bound on
the A100, which raises the memory ceiling without raising per-SM decode
throughput: they cross the ridge point of the roofline~\citep{roofline2009}. The floor crosses too --- a kernel reading no
weights is worth $4.77\times$ FP16 on the L40S and $1.68\times$ on the A100,
shrinking what our launch geometry buys from nearly five to under two. The
4-bit competitor crosses from further away ($3.38\times$ to
$1.82\times$), placing the effect in the whole class of decoding kernels,
lattice or otherwise.

What survives is the \emph{ordering}: \textsc{Planes14} leads, the two
\textsc{Golay70} decoders trail, v2 beats v1 --- the format work
is not an artifact of one machine --- and the f64 check holds to the last
digit. The headline is bounded: \emph{decoding at matvec speed is a result
about a bandwidth-bound regime}, and this point measures where it ends. On an
HBM-class card, expect Table~\ref{tab:layouts}'s ordering and none of its
ratios.
